% GENERATED FILE. Edit canonical lessons, then run npm run generate:books.
% canonical-source-hash: fnv1a64:eb4465b632d02ac2
% canonical-lessons: ML-C164-aviyal, ML-C164-pappatam, ML-C164-cutulla, ML-C164-tanuppulla, ML-C164-nilamulla

\chapter{Aviyal, Pappadam, Hot, Cold, Long}
\label{ch:ml-aviyal-pappadam-hot-cold-long}
\hlchaptermodality{164}

\begin{chapteropening}
\textbf{By the end of this chapter:} \emph{I can say aviyal, a pappadam, hot, cold and long.}

Complete the last lesson of chapter 164: I can say aviyal, a pappadam, hot, cold and long.
\end{chapteropening}

\section[\texorpdfstring{\ml{അവിയൽ} (aviyal)}{അവിയൽ (aviyal)}]{\ml{അവിയൽ} (aviyal) — aviyal, a mixed-vegetable dish}

\label{lesson:ML-C164-aviyal}



\begin{quote}
Before the new one: say the Malayalam for beans, then the Malayalam for dal.

Ask \textquotedblleft{}where from?\textquotedblright{} (\emph{Eviṭe ninnŭ?})
\end{quote}

\subsection*{You'll want to know: \ml{അവിയൽ}}
\textbf{\ml{അവിയൽ}} — \emph{aviyal} — \textquotedblleft{}aviyal, a mixed-vegetable dish\textquotedblright{}.

\begin{cousinweb}
\textbf{\ml{അവിയൽ}} (\emph{aviyal}) is native Dravidian, from the verb \textquotedblleft{}to boil, to cook in steam\textquotedblright{}; DEDR 268 lists the root with Tamil \emph{avi}.
\end{cousinweb}

\begin{grammarlens}[title={What you've built}]
One more word for this chapter.
\end{grammarlens}

\subsection*{Your turn}
\begin{itemize}
\raggedright
  \item \emph{Say it:} \emph{aviyal}
  \item \emph{Say it:} \emph{aviyal}, once more
  \item \emph{Say it:} say \emph{aviyal}
  \item \emph{Recall:} say the Malayalam for beans, then the Malayalam for dal, then say \emph{aviyal} again
\end{itemize}

\begin{tcolorbox}[breakable,colback=teal!4,colframe=teal!35!black,title={Before you move on}]
What does \ml{അവിയൽ} mean? (\textquotedblleft{}Aviyal, a mixed-vegetable dish\textquotedblright{}.) Say it once more.
\end{tcolorbox}
\section[\texorpdfstring{\ml{പപ്പടം} (pappaṭaṁ)}{പപ്പടം (pappaṭaṁ)}]{\ml{പപ്പടം} (pappaṭaṁ) — a pappadam}

\label{lesson:ML-C164-pappatam}



\begin{quote}
Before the new one: say the Malayalam for dal, then the Malayalam for aviyal.
\end{quote}

\subsection*{You'll want to know: \ml{പപ്പടം}}
\textbf{\ml{പപ്പടം}} — \emph{pappaṭaṁ} — \textquotedblleft{}a pappadam\textquotedblright{}.

\begin{cousinweb}
\textbf{\ml{പപ്പടം}} (\emph{pappaṭaṁ}) is native Dravidian, one word with Tamil \emph{pappaṭam}, Kannada \emph{pappaḍi} and Telugu \emph{appaḍamu} (DEDR 3928). DEDR 3928 compares Sanskrit \emph{parpaṭa}.
\end{cousinweb}

\begin{grammarlens}[title={What you've built}]
One more word for this chapter.
\end{grammarlens}

\subsection*{Your turn}
\begin{itemize}
\raggedright
  \item \emph{Say it:} \emph{pappaṭaṁ}
  \item \emph{Say it:} \emph{pappaṭaṁ}, once more
  \item \emph{Say it:} say \emph{pappaṭaṁ}
  \item \emph{Recall:} say the Malayalam for dal, then the Malayalam for aviyal, then say \emph{pappaṭaṁ} again
\end{itemize}

\begin{tcolorbox}[breakable,colback=teal!4,colframe=teal!35!black,title={Before you move on}]
What does \ml{പപ്പടം} mean? (\textquotedblleft{}A pappadam\textquotedblright{}.) Say it once more.
\end{tcolorbox}
\section[\texorpdfstring{\ml{ചൂടുള്ള} (cūṭuḷḷa)}{ചൂടുള്ള (cūṭuḷḷa)}]{\ml{ചൂടുള്ള} (cūṭuḷḷa) — hot}

\label{lesson:ML-C164-cutulla}



\begin{quote}
Before the new one: say the Malayalam for aviyal, then the Malayalam for a pappadam.

Say ninety and a hundred. (\emph{Toṇṇūṟŭ}, \emph{nūṟŭ}.)
\end{quote}

\subsection*{You'll want to know: \ml{ചൂടുള്ള}}
\textbf{\ml{ചൂടുള്ള}} — \emph{cūṭuḷḷa} — \textquotedblleft{}hot\textquotedblright{}.

\begin{cousinweb}
\textbf{\ml{ചൂടുള്ള}} (\emph{cūṭuḷḷa}) is native Dravidian, \textquotedblleft{}having heat\textquotedblright{}: \emph{cūṭŭ}, \textquotedblleft{}heat\textquotedblright{} (DEDR 2654), with the ending \emph{-uḷḷa}, \textquotedblleft{}that has\textquotedblright{}, from the verb \textquotedblleft{}to be\textquotedblright{} (DEDR 697, with Tamil \emph{uḷ}). DEDR 2654 compares Sanskrit \emph{cūḍ}.
\end{cousinweb}

\begin{grammarlens}[title={What you've built}]
One more word for this chapter.
\end{grammarlens}

\subsection*{Your turn}
\begin{itemize}
\raggedright
  \item \emph{Say it:} \emph{cūṭuḷḷa}
  \item \emph{Say it:} \emph{cūṭuḷḷa}, once more
  \item \emph{Say it:} say \emph{cūṭuḷḷa}
  \item \emph{Recall:} say the Malayalam for aviyal, then the Malayalam for a pappadam, then say \emph{cūṭuḷḷa} again
\end{itemize}

\begin{tcolorbox}[breakable,colback=teal!4,colframe=teal!35!black,title={Before you move on}]
What does \ml{ചൂടുള്ള} mean? (\textquotedblleft{}Hot\textquotedblright{}.) Say it once more.
\end{tcolorbox}
\section[\texorpdfstring{\ml{തണുപ്പുള്ള} (taṇuppuḷḷa)}{തണുപ്പുള്ള (taṇuppuḷḷa)}]{\ml{തണുപ്പുള്ള} (taṇuppuḷḷa) — cold}

\label{lesson:ML-C164-tanuppulla}



\begin{quote}
Before the new one: say the Malayalam for a pappadam, then the Malayalam for hot.
\end{quote}

\subsection*{You'll want to know: \ml{തണുപ്പുള്ള}}
\textbf{\ml{തണുപ്പുള്ള}} — \emph{taṇuppuḷḷa} — \textquotedblleft{}cold\textquotedblright{}.

\begin{cousinweb}
\textbf{\ml{തണുപ്പുള്ള}} (\emph{taṇuppuḷḷa}) is native Dravidian, \textquotedblleft{}having cold\textquotedblright{}: \emph{taṇuppŭ}, \textquotedblleft{}coldness\textquotedblright{} (DEDR 3045), with \emph{-uḷḷa}, \textquotedblleft{}that has\textquotedblright{}, from the verb \textquotedblleft{}to be\textquotedblright{} (DEDR 697).
\end{cousinweb}

\begin{grammarlens}[title={What you've built}]
One more word for this chapter.
\end{grammarlens}

\subsection*{Your turn}
\begin{itemize}
\raggedright
  \item \emph{Say it:} \emph{taṇuppuḷḷa}
  \item \emph{Say it:} \emph{taṇuppuḷḷa}, once more
  \item \emph{Say it:} say \emph{taṇuppuḷḷa}
  \item \emph{Recall:} say the Malayalam for a pappadam, then the Malayalam for hot, then say \emph{taṇuppuḷḷa} again
\end{itemize}

\begin{tcolorbox}[breakable,colback=teal!4,colframe=teal!35!black,title={Before you move on}]
What does \ml{തണുപ്പുള്ള} mean? (\textquotedblleft{}Cold\textquotedblright{}.) Say it once more.
\end{tcolorbox}
\section[\texorpdfstring{\ml{നീളമുള്ള} (nīḷamuḷḷa)}{നീളമുള്ള (nīḷamuḷḷa)}]{\ml{നീളമുള്ള} (nīḷamuḷḷa) — long}

\label{lesson:ML-C164-nilamulla}



\begin{quote}
Before the new one: say the Malayalam for hot, then the Malayalam for cold.
\end{quote}

\subsection*{You'll want to know: \ml{നീളമുള്ള}}
\textbf{\ml{നീളമുള്ള}} — \emph{nīḷamuḷḷa} — \textquotedblleft{}long\textquotedblright{}.

\begin{cousinweb}
\textbf{\ml{നീളമുള്ള}} (\emph{nīḷamuḷḷa}) is native Dravidian, \textquotedblleft{}having length\textquotedblright{}: the noun for \textquotedblleft{}length, distance\textquotedblright{} (DEDR 3692, with Tamil \emph{nīḷ}, \textquotedblleft{}to grow long\textquotedblright{}), with \emph{-uḷḷa}, \textquotedblleft{}that has\textquotedblright{} (DEDR 697).
\end{cousinweb}

\begin{grammarlens}[title={What you've built}]
One more word for this chapter.
\end{grammarlens}

\subsection*{Your turn}
\begin{itemize}
\raggedright
  \item \emph{Say it:} \emph{nīḷamuḷḷa}
  \item \emph{Say it:} \emph{nīḷamuḷḷa}, once more
  \item \emph{Say it:} say \emph{nīḷamuḷḷa}
  \item \emph{Recall:} say the Malayalam for hot, then the Malayalam for cold, then say \emph{nīḷamuḷḷa} again
\end{itemize}

\begin{tcolorbox}[breakable,colback=teal!4,colframe=teal!35!black,title={Before you move on}]
What does \ml{നീളമുള്ള} mean? (\textquotedblleft{}Long\textquotedblright{}.) Say it once more.
\end{tcolorbox}
